\documentclass[11pt,a4paper]{article}

\usepackage[utf8]{inputenc}
\usepackage[T1]{fontenc}
\usepackage{amsmath,amssymb}
\usepackage{geometry}
\geometry{margin=0.85in}
\usepackage{hyperref}
\hypersetup{colorlinks=true,linkcolor=blue,citecolor=blue,urlcolor=blue}
\usepackage{booktabs}
\usepackage{array}
\usepackage{tabularx}
\usepackage{enumitem}
\usepackage{caption}
\usepackage{multirow}
\usepackage{xcolor}
\usepackage{fancyhdr}

\setlist{nosep,leftmargin=1.2em}
\setlength{\parskip}{0.4em}
\setlength{\parindent}{0em}

\pagestyle{fancy}
\fancyhf{}
\fancyhead[L]{\small V2 Experiment Suite — Cahn-Hilliard PINN}
\fancyhead[R]{\small DRP, Brown University}
\fancyfoot[C]{\thepage}

\title{\vspace{-1.5cm}\textbf{Experiment Suite v2: Quasi-Newton Optimizers\\for the Cahn-Hilliard PINN}\vspace{-0.3cm}}
\author{Maria Macías Ordóñez \\ \small Directed Reading Program, Brown University}
\date{March 2026}

\begin{document}
\maketitle
\thispagestyle{fancy}

% ==========================================================================
\section{Motivation}
% ==========================================================================

Kiyani et al.~\cite{kiyani2025} demonstrated that self-scaled quasi-Newton methods (SSBroyden) outperform standard BFGS and L-BFGS by 2--3 orders of magnitude on five benchmark PDEs (Burgers, Allen-Cahn, Kuramoto-Sivashinsky, Ginzburg-Landau, Stokes).
Their study does \emph{not} include the Cahn-Hilliard (CH) equation, which introduces two challenges absent from those benchmarks: (i)~a \textbf{biharmonic operator} $\nabla^4 U$ producing a fourth-order spatial stiffness that severely conditions the loss Hessian, and (ii)~a \textbf{conservation law} $\int U\,d\Omega = \text{const}$ that PINNs must satisfy implicitly.

This experiment suite extends the paper's methodology to the CH equation:
\begin{equation}
    \frac{\partial U}{\partial t} \;=\; D\,\nabla^2\!\bigl(\nabla^2 U + a_2\,U + a_4\,U^3\bigr),
    \qquad (x,y)\in[0,128]^2,\; t\in[0,20],
    \label{eq:ch}
\end{equation}
with $D=a_2=a_4=1$ and periodic boundary conditions.
We replicate all ablation axes from~\cite{kiyani2025} and add one novel axis (RAdam warmup) to test whether rectified adaptive learning rates improve the quality of the first-order preconditioning phase.

% ==========================================================================
\section{Canonical Configuration}
% ==========================================================================

Every experiment inherits from a single canonical template and modifies exactly one ablation axis.
The canonical matches the paper's Burgers Case~1 budget, which achieved $\ell_2 \sim 10^{-8}$ with SSBroyden~\cite[Table~2]{kiyani2025}.

\begin{center}
\small
\begin{tabular}{@{}ll@{\qquad}ll@{}}
\toprule
\textbf{Parameter} & \textbf{Value} & \textbf{Parameter} & \textbf{Value} \\
\midrule
Network        & $4\times20$ (1,361 params) & Adam warmup  & 1,000 epochs \\
Activation     & $\tanh$                    & Adam LR      & $10^{-3}$, decay 0.98 \\
BFGS variant   & SSBroyden2                 & BFGS iters   & 50,000 \\
$N_\text{change}$ & 500                     & $N_\text{int}$ & 10,000 \\
$N_\text{IC}$  & 1,000                      & $N_\text{BC}$  & 500 \\
RAD            & enabled ($k_1\!=\!k_2\!=\!1$) & Loss weights & $\lambda_\text{pde}\!=\!1,\;\lambda_\text{ic}\!=\!\lambda_\text{bc}\!=\!5$ \\
Precision      & \texttt{float64}           & Seeds         & 1, 2, 3 \\
\bottomrule
\end{tabular}
\end{center}

% ==========================================================================
\section{Experiment Suite}
% ==========================================================================

The suite comprises \textbf{20 configurations} organized in 7 groups (A--G), each run with 3 random seeds for a total of \textbf{60 jobs} on Brown's Oscar cluster (48\,h wall-time, 4\,GB RAM, CPU-only \texttt{float64}).
Table~\ref{tab:suite} lists every experiment with its single ablation change relative to canonical.

\begin{table}[!ht]
\centering
\caption{Complete experiment suite. Each row overrides exactly one axis from the canonical configuration (Section~2). ``Paper ref.''\ indicates the corresponding table, figure, or case in~\cite{kiyani2025}.}
\label{tab:suite}

\vspace{0.3em}
\small
\setlength{\tabcolsep}{4pt}
\begin{tabularx}{\textwidth}{@{}l l l >{\raggedright\arraybackslash}X l@{}}
\toprule
\textbf{ID} & \textbf{Name} & \textbf{Change from canonical} & \textbf{Rationale} & \textbf{Paper ref.} \\
\midrule
\multicolumn{5}{@{}l}{\textit{Group A — Optimizer Variant Comparison (core result)}} \\[2pt]
A1 & BFGS          & variant = BFGS         & Textbook quasi-Newton baseline; no self-scaling.                & Tab.~2, Fig.~3 \\
A2 & SSBFGS-OL     & variant = SSBFGS\_OL   & Oren-Luenberger scaling $\tau_k=1/b_k$.                       & Tab.~2 \\
A3 & SSBFGS-AB     & variant = SSBFGS\_AB   & Al-Baali scaling $\tau_k=\min(1,1/b_k)$; clips over-scaling.  & Tab.~2 \\
A4 & SSBroyden1    & variant = SSBroyden1   & Broyden class ($\phi\!\to\!1$, DFP direction); DRP-patched.    & Tab.~2, Fig.~5 \\
A5 & SSBroyden2    & \textbf{canonical}     & Broyden class ($\phi\!\to\!0$, BFGS direction); best in paper. & Tab.~2, Fig.~5 \\
A6 & L-BFGS        & method = L-BFGS-B      & Limited-memory; $O(mN)$ storage; worst quasi-Newton in paper.  & Tab.~2, Fig.~4 \\
\midrule
\multicolumn{5}{@{}l}{\textit{Group B — Adam Warmup Duration}} \\[2pt]
B1 & Warmup 0      & Adam epochs = 0        & No first-order preconditioning; cold start for BFGS.           & Case~3 \\
B2 & Warmup 5K     & Adam epochs = 5,000    & Matches Allen-Cahn warmup; deeper loss basin entry.            & Tab.~7 \\
B3 & Warmup 10K    & Adam epochs = 10,000   & Tests diminishing returns of extended warmup.                  & --- \\
\midrule
\multicolumn{5}{@{}l}{\textit{Group C — Network Architecture}} \\[2pt]
C1 & 8$\times$20   & layers = [3,20$^8$,1]  & Deeper (3,041 params); parallels Burgers Case~4.               & Case~4 \\
C2 & 5$\times$30   & layers = [3,30$^5$,1]  & Wider (3,871 params); parallels GL Case~1 architecture.        & Tab.~10 \\
\midrule
\multicolumn{5}{@{}l}{\textit{Group D — First-Order Baseline}} \\[2pt]
D1 & Adam only     & 51K Adam, 0 BFGS       & Same iteration budget; establishes the performance floor.       & All tables \\
\midrule
\multicolumn{5}{@{}l}{\textit{Group E — RAdam Warmup (DRP extension)}} \\[2pt]
E1 & RAdam 1K      & optimizer = RAdam      & Rectified Adam~\cite{liu2020radam}; built-in LR warmup.        & \textit{novel} \\
E2 & RAdam 5K      & RAdam + 5K epochs      & Longer RAdam; cf.\ B2 (Adam 5K) for controlled comparison.     & \textit{novel} \\
\midrule
\multicolumn{5}{@{}l}{\textit{Group F — Collocation \& Resampling}} \\[2pt]
F1 & $N_\text{chg}$=200  & batch\_size = 200       & Frequent resampling; best in v1 ($\ell_2\!=\!0.85$).    & Sec.~3.1 \\
F2 & $N_\text{chg}$=1000 & batch\_size = 1,000     & Infrequent resampling; risk of overfitting fixed set.    & Sec.~3.1 \\
F3 & 20K points    & $N_\text{int}$=20K, $N_\text{IC}$=2K & Double collocation; better PDE residual coverage. & Sec.~3.1 \\
\midrule
\multicolumn{5}{@{}l}{\textit{Group G — Loss Conditioning}} \\[2pt]
G1 & $H_0$ scaling & initial\_scale = true   & Nocedal-Wright first-step Hessian scaling.                     & Sec.~2.3 \\
G2 & Power $p\!=\!2$     & power = 2              & Minimizes $\mathcal{L}^{1/2}$; flattens loss landscape.  & Sec.~2.4 \\
G3 & Power $p\!=\!4$     & power = 4              & Stronger flattening; second-best $\ell_2$ in v1.         & Sec.~2.4 \\
\bottomrule
\end{tabularx}
\end{table}

% ==========================================================================
\section{Ablation Rationale}
% ==========================================================================

\textbf{Group A} replicates the paper's central finding: on Burgers' equation, SSBroyden achieves $\ell_2 \sim 10^{-8}$ while standard BFGS reaches only ${\sim}10^{-5}$ and L-BFGS stalls at ${\sim}10^{-3}$~\cite[Table~2]{kiyani2025}.
The self-scaling factor $\tau_k$ rescales the inverse Hessian approximation at every iteration, and the Broyden-class formulation optimally interpolates between BFGS and DFP updates via a parameter $\phi_k$.
Extending this to the CH equation tests whether these gains persist under fourth-order stiffness.

\textbf{Group B} follows the paper's observation that Adam warmup duration affects quasi-Newton convergence.
Burgers Cases 1--3 compare 1K vs.\ 0 warmup~\cite[Table~2]{kiyani2025}; Allen-Cahn uses 5K~\cite[Table~7]{kiyani2025}.
The CH biharmonic operator may need longer warmup due to its higher condition number.

\textbf{Group C} addresses a practical finding from our v1 suite: all 17 experiments produced $\ell_2 \sim O(1)$ against a spectral reference solution, suggesting the $4\times20$ network (1,361 parameters) lacks capacity for the $128\times128$ spatial domain.
The paper uses $8\times20$ for Burgers Case~4 and $5\times30$ for Ginzburg-Landau~\cite[Tables~2,~10]{kiyani2025}.

\textbf{Group D} provides the first-order baseline. The paper shows Adam alone consistently stalls orders of magnitude above quasi-Newton final losses across all PDEs.

\textbf{Group E} is our \textbf{novel contribution}. RAdam~\cite{liu2020radam} rectifies adaptive learning rate variance during early training, acting as automatic warmup. The paper exclusively uses standard Adam; we test whether RAdam produces a better initialization for the BFGS phase.

\textbf{Group F} tests collocation strategy. The paper uses residual-based adaptive distribution (RAD) with periodic resampling every $N_\text{change}$ iterations~\cite[Sec.~3.1]{kiyani2025}. Our v1 results found $N_\text{change}=200$ outperformed 500 and 1000, and F3 tests whether under-resolution of the $\nabla^4$ residual is a bottleneck.

\textbf{Group G} tests two loss conditioning techniques: Nocedal-Wright initial scaling~\cite[Sec.~2.3]{kiyani2025} and the power transform $\tilde{\mathcal{L}} = \mathcal{L}^{1/p}$~\cite[Sec.~2.4]{kiyani2025}, which flattens the loss surface near sharp minima.

% ==========================================================================
\section{Evaluation}
% ==========================================================================

Each experiment is evaluated against a spectral reference solution (ETDRK4, $128\times128$, $\Delta t = 10^{-3}$):
\textbf{(1)}~Relative $\ell_2$ error $e = \|U_\text{PINN} - U_\text{ref}\|_2 / \|U_\text{ref}\|_2$ on 101 snapshots;
\textbf{(2)}~mass conservation $\Delta M = \max_t |\langle U(t)\rangle - \langle U(0)\rangle|$;
\textbf{(3)}~final loss and per-component losses;
\textbf{(4)}~wall-clock time.
All metrics are mean $\pm$ std over 3 seeds.
The 60 jobs run as a single SLURM array on Oscar (48\,h, 4\,GB, 4 CPUs, no GPU).

% ==========================================================================
\begin{thebibliography}{9}
\bibitem{kiyani2025}
R.~Kiyani, O.~Manga, S.~Zeng, and G.~E.~Karniadakis,
``Optimizing the Optimizer for Data-Free Physics-Informed Neural Networks and Kolmogorov-Arnold Networks,''
\textit{Comput.\ Methods Appl.\ Mech.\ Engrg.}, vol.~446, 118308, 2025.
\href{https://arxiv.org/abs/2501.16371}{arXiv:2501.16371v5}.

\bibitem{liu2020radam}
L.~Liu, H.~Jiang, P.~He, W.~Chen, X.~Liu, J.~Gao, and~J.~Han,
``On the Variance of the Adaptive Learning Rate and Beyond,''
\textit{Proc.\ ICLR}, 2020.
\end{thebibliography}

\end{document}
